% ============================================
% Discussion
% ============================================
\section{Discussion}
\label{sec:discussion}


Our work demonstrates that effective interaction intensity, environmental feedbacks, and assembly-generated interaction structure govern community coalescence outcomes, determining both the dominant outcome type and the level at which selection operates. In both experiments and modeling, when effective interactions are strong, Dominance prevails and species persistence is correlated within parental communities, indicating community-level selection; when interactions are weak, Mixture predominates and species respond independently. \rev{Furthermore, predictability from dominant species reveals two distinct routes associated with Dominance: a top-down-associated regime in which dominant taxa and environmental feedbacks, including pH modification, may bias winner identity, and an emergent regime where collective dynamics reduce predictability.} These findings reconcile the Clements--Gleason debate: effective interaction intensity determines whether communities behave as cohesive units or loose species assemblages.


Our work provides the first experimental demonstration of alternative regimes in community coalescence, reconciling previously contrasting observations of community-level versus species-level selection. Prior studies have reported both outcomes: some found communities behaving as cohesive units\citep{Tikhonov2016, DiazColunga2022}, while others observed species responding independently\citep{VanderGucht2007, Vermeij1991}. Our results suggest that these conflicting observations may reflect differences in interaction strength across experimental systems. For instance, the absence of community-level selection in \textit{in vitro} gut microbiome coalescence experiments\citep{Goldman2025, Walton2025} is consistent with our framework. Rich media such as BHI are analogous to our Nutr$-$ condition, providing complex and diverse resources while lacking high concentrations of readily metabolizable carbon sources. This resource structure reduces both direct resource competition and pH-mediated interactions, thereby weakening interspecies interactions\citep{Ratzke2020} and favoring species-level rather than community-level dynamics. Likewise, interaction strength may vary systematically across ecosystems: microbial biofilms experience strong metabolic interactions due to high cell densities and shared resources\citep{Nadell2016}, whereas mobile macroscopic organisms interact more transiently across larger spatial scales. These differences may explain why community-level selection is more frequently observed in microbial coalescence studies. This framework provides a unified perspective: rather than asking whether communities are cohesive units, we should ask under what conditions they become so, with interaction strength emerging as a key determinant.


\rev{Notably, and perhaps counterintuitively, our results demonstrate that community-level cohesion can emerge without cooperative interactions between species, a prediction made by Tikhonov \citep{Tikhonov2016} using resource-competition models and consonant with classical random Lotka--Volterra theory \citep{May1972, Grilli2017}. In our system, purely competitive interactions, when structured by assembly history, are sufficient to produce the correlated species fates that underlie Dominance. This ``cohesion without cooperation'' arises because assembly filters species into mutually weakly competing groups whose members face stronger competition from outsiders.}


A growing body of work has identified interaction strength as a key control parameter governing diverse aspects of microbial community dynamics, including diversity \citep{Hu2022, Marsland2019}, stability \citep{Allesina2012, Ratzke2020}, coexistence \citep{Grilli2017}, and invasion outcomes \citep{Hu2025, Kurkjian2021}. Our results extend this claim to include coalescence, providing further evidence that interaction strength determines when communities exhibit collective behavior. The parallel with priority effects is particularly instructive; recent work showed that assembly order shapes final composition only under strong interactions, whereas weak interactions lead to convergent outcomes regardless of assembly history \citep{Hu2025, Fukami2015}. Our coalescence results mirror this pattern: strong interactions preserve parental identity, while weak interactions yield convergent Mixtures. Together, these findings reinforce interaction strength as a coarse-grained parameter that predicts history-dependent community dynamics across multiple ecological contexts. \rev{Importantly, this parameter is phenomenological: it describes effective species coupling in the normalized community dynamics and can incorporate multiple underlying ecological mechanisms.}

\rev{We define community-level selection operationally as origin-correlated persistence during coalescence, rather than as a claim about one exclusive biochemical mechanism \citep{Mansour2018,Rillig2015}. This is a phenomenological description of the outcome, not a mechanistic attribution on its own. Accordingly, our evidence for community-level selection is operational and outcome-based: we assess it from parental-community similarity after coalescence and pairwise species-fate correlation, rather than directly measuring the mechanistic selection pressures that generate these outcomes. Under this definition, shared environmental tolerances, correlated traits, and environmental filtering by dominant pH-modifying taxa are alternative mechanistic routes that can produce the same outcome-level pattern when they make parental-community origin predictive of persistence. Several controls narrow purely passive explanations. The positive pairwise selection correlation is not driven solely by the most abundant species (Supplementary Information), arguing against simple hitchhiking. The assembly-history comparison, in which pre-assembled communities show higher Dominance than directly assembled ones (Extended Data Fig.~7, Supplementary Fig.~19), shows that shared ecological history contributes to correlated persistence. The gLV model, which lacks explicit environmental filtering and assigns identical growth rates to all species, further shows that pairwise competitive structure is sufficient to generate Dominance and positive within-community selection correlation. Thus, our scope here is to define and test community-level selection as an outcome-level pattern, while future work can directly disentangle how different mechanisms contribute across coalescence regimes.}

Our experimental design uses nutrient concentration to perturb net interaction intensity and environmental feedbacks, but this should not be read as a direct monotonic mapping from external resource supply to effective pairwise Lotka--Volterra coefficients. The gLV parameter $\mu$ is a simplified scalar model parameter, and the mapping between nutrient concentration and $\mu$ is necessarily approximate. This phenomenological mapping is intended to capture the net interaction intensity expressed in invasion outcomes, rather than to identify a unique biochemical mechanism. In our experiments, the nutrient-dependent increase in invasion resistance could reflect several non-mutually exclusive processes, including denser resource competition, increased metabolic activity, environmental modification such as pH shifts, changes in carrying capacities, and changes in competition--facilitation balance. This broader interpretation is consistent with work showing that microbial communities can assemble through environmental feedbacks and nutrient-mediated ecological structure \citep{Goldford2018,Estrela2021}, and with consumer-resource theory showing that external supply rate need not map monotonically onto effective Lotka--Volterra coefficients under broad conditions \citep{DuanPawar2025}. Nevertheless, the monotonic increase in failed pairwise invasions with nutrient supply provides direct experimental evidence that invasion resistance increases across our nutrient gradient, and the qualitative agreement between model predictions and experimental outcomes across all three conditions supports the utility of this coarse-grained framework.

\rev{A related caveat concerns our natural community experiments. The stabilization phase, seven serial growth-dilution cycles in defined laboratory media, may pre-select natural communities toward taxa and metabolic strategies that thrive under these culture conditions \citep{Goldford2018}, potentially reducing effective ecological heterogeneity and making natural sample-derived communities more similar to our synthetic consortia. Our post-stabilization 16S data show that this filtering was not complete at the ASV level: natural sample-derived parental communities retained higher richness than the synthetic parental communities ($13.7 \pm 7.2$ versus $9.8 \pm 4.8$ ASVs above the 0.1\% threshold) and low ASV overlap among communities from different source samples (Supplementary Figs.~22--25). However, because we did not sequence the original environmental inocula before enrichment, we cannot quantify taxonomic convergence during stabilization. We also cannot directly test functional convergence because these experiments used 16S community profiling rather than functional metagenomic, metabolomic, or trait-resolved measurements. Thus, pre-selection could contribute to the convergent coalescence patterns observed between synthetic and natural sample-derived communities, and the natural-community results should be interpreted as evidence that the nutrient-dependent increase in Dominance occurs in taxonomically richer laboratory-stabilized communities, not as evidence for unrestricted generality across unfiltered natural ecosystems. Future work using culture-independent time series or broader environmental sampling could help disentangle laboratory adaptation from intrinsic ecological dynamics.}

In nature, coalescence occurs in richer settings that may shift regimes and outcomes. Environmental heterogeneity (temperature, pH buffering, and other physicochemical factors) covaries with interaction strength and can reshape coalescence across habitats \citep{Rillig2015, Castledine2020, Rillig2017, Tropini2017}. In host-associated microbiomes, host filtering, priority effects, and spatial structure further influence successful colonization, making coalescence a useful framework for predicting and steering microbiome transfer \citep{Liu2024, Smillie2018, Rocca2021}. Our experiments focus on steady-state outcomes and thus do not capture temporal dynamics such as fluctuating resources, migration pulses, or host responses. \rev{Additionally, both our experimental system and theoretical model focus on a competition-dominated regime; systems with substantial mutualism or facilitation may exhibit qualitatively different dynamics.} Incorporating these axes into both theory and experiment will help build a more comprehensive picture of community-level selection across habitats and scales.
